\section{Introduction}
\label{sec:introduction}

The YOSO (You Only Speak Once) model was introduced as a way to obtain secure
multiparty computation against a very strong adaptive adversary
\cite{GentryEtAl2021YOSOFull}. A protocol role is assigned to an ephemeral
party which prepares one message, deletes the secret state used to prepare
it, and then sends the message. Only at this point does the role reveal
itself. Corrupting the sender after publication is therefore too late: the
message is already available and the secrets that produced it no longer
exist. Later work has developed this anonymous, ephemeral-role paradigm in
several directions; the erase-before-send discipline remains central.

This paper starts from the deliberately unfriendly question: what if deletion
is impossible? The restriction might be legal---a regime may require state
to be retained for later inspection---or technological. For example, a
party running in a cloud environment might fear that speaking identifies its
machine and enables the cloud operator to inspect a residual snapshot despite
an attempted deletion. We capture the extreme version of this concern by
declaring that whenever an ordinary party speaks, its complete state leaks to
the adversary.

The immediate answer seems to be ``obviously impossible.'' Eventually all
parties relevant to the computation will speak and hence reveal everything.
This conclusion overlooks one asymmetry. Speaking is optional. A machine may
remain silent forever, in which case its secret state may stay hidden. We
therefore place mutually exclusive pieces of state behind different
public-key roles. The selected roles speak and are expended; the unselected
roles never speak. In this sense, silence replaces erasure as the one-time
resource.

We call the resulting setting \emph{YAMSO}: You At Most Speak Once. In a
YAMSO execution, a synchronous ledger maintains a growing ordered list of
public keys and messages. The public-key index determines the protocol role.
For the preliminary model, each newly registered computation role is corrupted
independently with probability $p$, whereas named application parties retain
the ordinary corruption interface. An honest role keeps its secret key inside
a protected ledger interface until its scheduled opportunity to speak. Once a
computation role speaks, its key and complete role state may become public;
its message is already irrevocable and the role is never used again.

The intended real-world interpretation has a much larger population of
apparently interchangeable machines. The adversary compromises a fixed
fraction of that population, while fresh clients and role states are placed
on machines randomly and without exposing the placement before a machine is
corrupted or speaks. Random placement then turns the fixed compromised set
into an approximately random intersection with each protocol committee. The
Bernoulli experiment in this paper deliberately collapses that mechanism into
a simple model; it does not model mobile corruption.

\paragraph{Finite feasibility.}
The basic two-phase blueprint isolates why the model need not be vacuous. For
each garbled input wire, the two structural labels are placed behind disjoint
committees, their semantic order is hidden, and an activation token causes
only one committee to speak. Its complete exposed state contains shares of a
label that is becoming public anyway. Honest members of the sibling committee
remain silent forever and protect the unused label. A missing or malformed
input closes onto a fixed structural branch whose hidden semantic value is a
uniform default bit.

\paragraph{Ongoing computation.}
The difficult part is not evaluating another garbled circuit; it is preserving
the hidden correlation between the current state-label pair and every later
stretch without giving a speaking linker both labels. We propose a finite
public \emph{oblivious garbled-circuit sampler} (OGCS), implemented
conjecturally by an obfuscated streaming program. It fixes one public
transition circuit, receives a growing sequence of certified public inputs,
and emits exactly prefix-preserving garbled stretches. The output-state label
pairs of one stretch are literally the input-state pairs of the next.

A copyable obfuscated program can otherwise be evaluated on invented public
keys and forked ledger histories. Our second use of silence is therefore a
one-time common-input certificate. Each bit of a high-distance encoding is
represented by two one-way preimages placed in different one-use roles. The
role matching the canonical message speaks; its opposite remains silent.
After the correct certificate is public, the speakers may reveal everything
without enabling a conflicting certificate. This observation removes the
long-lived certification key we initially appeared to need. The ledger still
provides a common immutable view, but it need not hold a signing secret.

This is an instance of the model changing the design intuition. If a single
certifier speaks, its leaked signing key destroys future authenticity. The
natural reaction is to declare certification impossible. YAMSO instead asks
which secret holders need never speak. The resulting construction is closely
related to known high-distance Lamport signatures, but its use as persistent
common truth under total post-speech exposure is the mechanism we need here.

\paragraph{Status and proof strategy.}
The ongoing construction is intentionally presented as a candidate, not a
theorem. We modularize its analysis by three whole-process replacement lemmas:
canonical one-use certification, programmable prefix-preserving OGCS, and
two-stage finite garbling/readout with full speaker exposure. This avoids a
naive UC-hybrid decomposition whose interfaces would export internal keys to
the environment. The first version fixes corrupt input providers before setup
and has public outputs. Two central realization gaps remain: the cited
general-purpose ordinary-iO-based succinct TM-iO is bounded-input rather
than horizon free, and predictable random-oracle points do not yet support
one online simulator. A hidden nonce prepared one
stretch ahead is the intended lift for the second gap.

\paragraph{Relation to prior work.}
Our starting point is YOSO \cite{GentryEtAl2021YOSOFull}. Anonymous blockchain
committees tolerating a corruption rate around one quarter are known, but
their secret handoff relies on erasures
\cite{BenhamoudaEtAl2020BlockchainSecret}. Blockchain one-time programs
already demonstrate conditional release of one garbled input-label vector on
an immutable ledger \cite{GoyalGoyal2017Blockchains}. No-erasure hidden
committees also appear in communication-local MPC, although there the selected
committee runs an ordinary multi-round computation
\cite{ChandranEtAl2024CommunicationLocalMPC}.

For certification, Katz and Vaikuntanathan give a direct high-distance
Lamport precedent \cite{KatzVaikuntanathan2009LeakageSignatures}, and Kelsey,
Lang, and Lucks study distributed hash-based signing and the
disjoint-honest-trustees one-time-key-reuse problem in the absence of
broadcast \cite{KelseyLangLucks2025DistributedHBS}. Neither result states
our one-shot common-input property under post-speech full exposure. For the
streaming component, relevant technical precedents include universal samplers,
cryptographic iterators, succinct and reusable garbling, TM obfuscation, and
streaming functional encryption
\cite{HofheinzEtAl2014UniversalSamplers,KoppulaLewkoWaters2014IOTM,%
LinPass2014SuccinctGarbling,GoldwasserEtAl2013TuringMachines,%
GuanKorbSahai2022StreamingFE}. We use these as component precedents rather
than claiming that an existing theorem already gives the YAMSO construction.
